\section{Methods}
\label{sec:methods}

Everything reported in Sections~\ref{sec:coverage} to~\ref{sec:forces} is a
single-point energy or a force at one fixed geometry. No structure moves during
any of those runs. The protocol removes every difference between the programs
being compared except the potential itself, so that a disagreement can only
come from the parameters or from the code that evaluates them. What a reader
needs only in order to rerun a job, the engine and reference control files,
the deck sizes, the \texttt{tleap} presets, the CHARMM36 route and the 1efs
coordinate edits, is tabulated in the Supporting Information and appears here
only where a choice carries an argument.

\subsection{The engine and how it was run}
\label{sec:methods-engine}

MOSAICS \cite{mosaics-software} is the natural-move Monte Carlo engine of Sim
\emph{et al.}\cite{sim2012naturalmoves} and Minary and
Levitt\cite{minary2010chainclosure}. It was built from the serial target,
\texttt{source/compile/serial}, of the engine's main development branch and
run from the command line. Every run directory in the evidence tree records the
SHA-256 of the binary that produced it and the commit it was built from, and
that is how a build is identified, because the banner the engine prints reads
\texttt{version.3.9.1} on every build used here. The released 3.9.1 binary was
compared against the development build on the same inputs before the
comparison began: it refuses the OL21/OL3, parmbsc1 and parmbsc0 decks at the
deck, because their torsion entries carry a per-molecule label it does not
know; it refuses the ff19SB input because of the correction-map term, and with
that term stripped it omits the term silently; and on the OL15/OL3, ff14SB and
CHARMM36 decks it runs and returns energies tens of thousands of \kcal{} from
the references, because it does not chain the residues it builds from a deck.
The development build reproduces every committed single point byte for byte.
The gate report is in the evidence tree. There is no software paper for the
engine, so it is cited by name and project page.

A MOSAICS job reads four things: a simulation input file, a residue topology
file (RTF) naming the atoms of each residue and their types, a set of parameter
files (a \emph{deck}), and a PDB structure. Supporting Information Table~S1
lists every setting that makes a run a single point. Four of them decide what
is measured. Every proposal width is zero, so the one Monte Carlo step taken
proposes the starting structure. \verb|\energy_report{2}| makes the engine
print each term and each part of each term, and \verb|\inter_list{none}|
switches the neighbour list off, with no interaction cut off or switched. The
nucleic runs use torsional proposals (\texttt{prop\_type} \texttt{tors}) with a
\texttt{time\_step\_md} of 0.4, the peptide runs Cartesian proposals
(\texttt{prop\_type} \texttt{cart}) with a time step of zero: in Cartesian mode a
non-zero step lets the structure drift even at zero proposal width, and the
energies then describe two different shapes. The fourteen
committed inputs of the two reference systems and the two controls also run
clean under the address sanitizer with \texttt{alloc\_dealloc\_mismatch} turned
on, which is off by default on macOS.

\paragraph{What is new in the engine, and what is not}
Most of this work is conversion and reference construction, not engine code.
The six energy terms (bond, bend, torsion and improper, one-four, non-bonded,
correction map), their analytic forces and the staging transform all predate
it. Three things changed in engine source. The deck reader carries a
per-molecule label on torsion entries and chains the residues it builds, which
is what the released 3.9.1 lacks. The correction-map analytic force was
identically zero at a dihedral of exactly $0$ or $\pi$, repaired in
\texttt{cmap.cpp} (finding F-21); the correction-map \emph{energy} was never
affected. And forty arrays allocated with \texttt{new[]} and grown with
\texttt{realloc} were made consistent across ten source files (F-19). The two
CHARMM terms AMBER does not have needed no engine change at all, riding on
records the engine already evaluates. Finding identifiers throughout follow
\texttt{evidence/findings.txt}, the register in which an identifier is never
reused.

\paragraph{The path a sampling run takes}
No run in Sections~\ref{sec:coverage} to~\ref{sec:forces} samples, so it is
worth saying which code a sampling run would execute. MOSAICS scores a proposal
by recomputing the whole potential: \texttt{Pos::potential}, in
\path{source/energy/energy.cpp}, calls the bond, bend, torsion, one-four,
non-bonded, correction-map and constraint routines over the entire molecule on
the trial coordinates, and the Metropolis ratio is formed from two such full
potentials. There is no incremental $\Delta E$ over a displaced subset, so the
term routines validated here are the term routines a proposal is judged by. One
difference remains. \verb|\inter_list| is \verb|none| in every single-point
run, so the non-bonded sum runs over all non-excluded pairs, while a production
run may select the link-list or neighbour-list branch; those branches restrict
the pair set of the same expression and are not exercised by any single point
here. The cost of a proposal is one full evaluation, quadratic in the number of
non-excluded pairs, and Section~\ref{sec:methods-cost} measures it.

\subsection{Force fields, and how a parameter set enters the engine}
\label{sec:methods-parameters}

MOSAICS reads a force field as a deck of five plain-text files alongside an
RTF: \path{.bond}, \path{.bend}, \path{.tors_and_impr}, \path{.onfo} for atoms
three bonds apart and \path{.vdw} for everything further. The decks
used here are \texttt{db\_ff14sb} and \texttt{db\_ff19sb} for protein;
\texttt{db\_bs0}, \texttt{db\_bsc1}, \texttt{db\_ol15\_ol3},
\texttt{db\_ol21\_ol3} and \texttt{db\_ol24\_ol3} for nucleic acid, the last
three each with a \texttt{\_terminal} variant; \texttt{db\_charmm36\_na}; and,
for the protein--RNA complex, two decks that hold protein and RNA types
together, \texttt{db\_ff14sb\_ol3} and \texttt{db\_ff19sb\_ol3}. Supporting
Information Table~S2 collects the entry counts the engine prints as it reads
the six decks of the two reference systems. How many of those entries a given
test system actually evaluates is not reported by the engine and is not
reported here.

The protein sets are ff14SB\cite{maier2015ff14sb} and
ff19SB\cite{tian2020ff19sb}. ff19SB replaces the ff14SB backbone dihedral terms
with per-residue two-dimensional $(\phi,\psi)$ correction maps, written in the
grid-based CMAP form introduced for
CHARMM\cite{mackerell2004cmap,mackerell2004cmapfull}, so its deck carries a
sixth file, \texttt{.cmap}, and its input carries \verb|\energy_term{cmap}| and
\verb|\cmap_database_file|. On the capped tetrapeptide the engine reads 44
database entries, loads 16 maps ($N=24$) and selects four of those sixteen, one
per interior residue.

The nucleic sets are parmbsc0\cite{perez2007bsc0}, parmbsc1\cite{ivani2016bsc1},
OL15\cite{zgarbova2015ol15}, OL21\cite{zgarbova2021ol21} and
OL24\cite{zgarbova2025ol24} for DNA, each paired with
$\chi_{\mathrm{OL3}}$\cite{zgarbova2011ol3} for RNA. Every torsion entry
carries a per-molecule label. The two bsc decks originally carried none, and
without it the four atom-type quadruplets that need different values for DNA
and for RNA collapse onto the DNA values, which is worth 0.47~\kcal{} on the
1efs torsion energy (F-03). The one-four electrostatic scale factor is stored
as exactly $5/6$ in every deck; storing it to seven or ten digits instead
shifts the 1efs total by $1.0\times10^{-4}$~\kcal{} (F-04).

Each deck is converted from a published parameter file and named after the
source it came from. The OL15/OL3, OL21/OL3 and ff14SB decks were written by
\texttt{tools/openmm\_xml\_to\_mosaics.py} from the OpenMM force-field XML
files distributed with OpenMM and openmmforcefields, overlaid with the exact
proper-torsion rows of a matched AMBER \texttt{tleap} topology so that wildcard
and exact torsion precedence match sander. The OL24/OL3 decks were written by
\texttt{tools/amber\_prmtop\_to\_mosaics.py} from resolved AMBER topologies,
the DNA half from the \texttt{ff-DNA-OL24} package distributed by the OL24
authors and the RNA half from the OL3 files in AmberTools. Regeneration was
checked with the unmodified converter as a control: \texttt{db\_ff14sb}
regenerates byte-identical; \texttt{db\_ol15\_ol3} does not regenerate at all,
and neither does the \texttt{tors\_and\_impr} file of \texttt{db\_ol21\_ol3}
(F-20), so what backs those two decks' provenance is a note in the repository
rather than a command a reader can run. The OL15/OL3 and OL21/OL3 rows should
be weighed accordingly.

\paragraph{Standard and terminal decks}
A standard nucleic deck expects a phosphate on every residue, which is what the
padded heteroduplex build provides. A chain with a real $5'$-hydroxyl end
selects different rows, so OL15, OL21 and OL24 each carry a
\texttt{\_terminal} deck written from true-terminal tetranucleotide topologies,
and it is the terminal deck that every free-ended system in this work reads.
parmbsc1 has no terminal deck. Run on the Drew--Dickerson dodecamer through its
standard deck it returns a torsion energy 0.385~\kcal{} above the references,
and the same experiment with OL21/OL3's standard deck in place of its terminal
one reproduces the size of that error (0.369~\kcal{}), so the gap is the missing
terminal profile and not the parameters. parmbsc0 does not run on that
structure at all. Neither combination is in the 26-cell matrix, and both are
recorded as open findings in the evidence tree.

\paragraph{CHARMM36}
The CHARMM36 nucleic acid parameters\cite{hart2012charmm36dna,denning2011charmm36rna}
reach the engine by a different route, recorded in
\texttt{evidence/charmm36.txt} and set out step by step in Supporting
Information section~S6. Deck and reference energies are both read out of an
OpenMM \texttt{System} built from one release, \texttt{charmm36\_2024.xml}, so
the comparison does not straddle two releases, but the numbers belong to that
one and no release-independent claim about CHARMM36 is made here. The two
CHARMM terms AMBER does not have ride on records the engine already has, the
Urey--Bradley $1$--$3$ harmonic on \texttt{\~{}bend\_parm} and the 188 harmonic
impropers on \texttt{\~{}torsion\_parm}, where the deck reader expands them into
a truncated sixth-order Fourier series. That series is the conversion's only
genuine approximation: exact to four figures at $10^{\circ}$ from equilibrium
and $-7.7$\,\% at $60^{\circ}$, with the intermediate angles in Table~S6, and
nucleobase impropers are stiff enough to stay inside $20^{\circ}$. CHARMM36
here is nucleic acid only, only on the $5'$-hydroxyl terminus, and only against
OpenMM; the protein force field\cite{best2012charmm36} is not implemented,
because two of its correction-map selectors distinguish a proline neighbour by
the class of the \emph{following} residue while \texttt{\~{}cmap\_binding}
selects a map by the residue itself, so supporting it changes the file format
and the engine (F-09).

\subsection{Reference engines}
\label{sec:methods-references}

Three independently written programs serve as references
(Table~\ref{tab:engines}), and where they agree with each other the target is
not in doubt. No number here is produced by a program that reads another
program's output; each writes its own output file and those files are
committed. Term-by-term cross-engine comparison of this kind has an established
precedent\cite{shirts2017engines}, and the unit and convention traps it has to
avoid are catalogued in ref.~\cite{braun2018bestpractices}.

\begin{table}[t]
\centering
\footnotesize
\setlength{\tabcolsep}{4pt}
\caption{The engines, the resolution of the file each number is read from, the
floor that resolution puts under a difference, and the committed file driving
each run. All four run in double precision. \texttt{D} is the directory for one
force field and system; ``d.p.'' is decimal places as printed. A composite term
sums two printed columns, so its floor is twice the per-column figure, about
$2\times10^{-4}$ for sander and $4.8\times10^{-7}$ for GROMACS.}
\label{tab:engines}
\begin{tabular}{@{}lllll@{}}
\toprule
Program & Version & Output resolution & Floor (\kcal{}) & Input \\
\midrule
MOSAICS & main branch, serial & 16 significant figures & not limiting & \texttt{D/mosaics/*.input} \\
sander  & 24.0 (AmberTools 24.8)\cite{case2023ambertools} & 4 d.p., \kcal{} & $10^{-4}$ & \texttt{D/sander/sander\_sp.in} \\
OpenMM  & 8.5.2\cite{eastman2017openmm7,eastman2024openmm8} & 12 d.p., \kcal{} & $5\times10^{-13}$ & \texttt{openmm\_single\_point.py} \\
GROMACS & 2025.4\cite{abraham2015gromacs} & 6 d.p., kJ/mol & $2.4\times10^{-7}$ & \texttt{D/gromacs/sp.mdp} \\
\bottomrule
\end{tabular}
\end{table}

Those floors are properties of the reference and not of MOSAICS, each a
half-unit in the last place that reference prints. OpenMM's is the one that
invites a misreading: it prints twelve decimal places in \kcal{}, not twelve
significant figures, so its floor is an absolute quantity and does not scale
with the size of the term. The rule this paper follows is that no difference
smaller than the coarser of the two engines' floors is reported as a
measurement; below that floor a difference is reported as below the resolution
of the reference output, with the floor stated, and is not read as agreement to
that many digits.

The three control files are reproduced option by option in Supporting
Information Table~S3. sander is run as a zero-cycle minimisation in vacuum with
no cutoff at dielectric 1, and its energies are read from the \texttt{FINAL
RESULTS} block. OpenMM is driven on the same topology and coordinate files
sander reads, on the Reference platform in double precision with no cutoff, no
constraints and no centre-of-mass removal. It reports one combined non-bonded
energy, so the system is evaluated twice more, once with every charge zeroed
and once with every Lennard-Jones $\epsilon$ zeroed, to split it; the two
halves and the combined value are all printed, and their difference is a
self-check the reader can confirm at $1.5\times10^{-11}$~\kcal{} on 1efs.
GROMACS is run as a rerun over the single frame with plain cut-off
electrostatics and Lennard-Jones, no modifiers, no dispersion correction, no
constraints, and the cutoff set past the far end of the molecule so that every
pair is inside it with 2\,nm to spare and no periodic image is; \texttt{sp.log}
confirms a potential shift of zero on every term. Energies are read from
\texttt{sp.xvg}, never from what \texttt{gmx energy} prints on screen, which is
rounded coarsely enough to look like a disagreement between engines. GROMACS
references exist for the six original AMBER combinations and for the ff14SB
proteins and the OL21/OL3, OL15/OL3 and OL3 nucleic systems; the ff19SB rows
carry no GROMACS total, for the reason given next.

\paragraph{The GROMACS conversion route, and what it does not carry}
The GROMACS topology is produced mechanically from the same AMBER topology the
other two engines read, by the ParmEd command in \texttt{evidence/engines.txt};
ParmEd is used as distributed with AmberTools\cite{case2023ambertools}. In the
topology this route produces, \texttt{[ cmaptypes ]} is keyed on the five
backbone atom types, which every residue of a protein shares, so ff19SB's
per-residue maps collapse onto one (F-06). On the four-residue peptide that
costs 0.30~\kcal{}; on ubiquitin the GROMACS correction-map energy is
61.54~\kcal{} against the 69.95 that MOSAICS and OpenMM both return. That is a
limitation of the conversion route as used here, not of the GROMACS file
format: published GROMACS ports of amino-acid-specific CMAP force fields
introduce residue-specific backbone atom types so that each map gets its own
type quintuple. That workaround was not attempted, so ff19SB is reported as
verified against sander and OpenMM, and the GROMACS column of every ff19SB row
is blank rather than wrong.

The engines do not group their terms identically, so all comparisons use three
composite terms. \emph{Electrostatics} means all of it: \texttt{EEL + 1-4 EEL}
in sander, \texttt{Coulomb-14 + Coulomb (SR)} in GROMACS, the epsilons-zeroed
evaluation in OpenMM, \texttt{Total Coulomb energy} in MOSAICS;
\emph{Lennard-Jones} is the matching sum, and \emph{dihedral} includes impropers
in all four programs.

\subsection{Test systems}
\label{sec:methods-systems}

Eleven structures are used (Table~\ref{tab:systems}, Figure~\ref{fig:systems}).
Two of them, the 818-atom DNA:RNA heteroduplex and the capped tetrapeptide,
carry the term-by-term comparison against all three references; the other nine
widen the chemistry the force fields are exercised on, from four nucleic bases
to eight, from four amino acids to twenty, and from interior residues to real
chain ends, disulfides, a bound peptide and a protein--RNA interface. A
single-point comparison is defined by its coordinates, so each file's provenance
is given here.

\begin{table}[t]
\centering
\footnotesize
\setlength{\tabcolsep}{3pt}
\caption{The eleven test systems. Atom counts are read from the built topology,
hydrogens included; the force fields are those with at least one reference
engine in the 26-cell matrix of Table~\ref{tab:matrix}. All are run in vacuum at
dielectric~1 with no cutoff. Source:
\texttt{work/preprint-I/systems/MANIFEST.json}, which also records the SHA-256 of
every coordinate file.}
\label{tab:systems}
% generated by make_tables.py from systems/MANIFEST.json and energies/MATRIX.json
\begin{tabular}{@{}llrrp{3.1cm}p{4.5cm}@{}}
\toprule
System & Class & Atoms & Chains & Force fields compared & What it tests \\
\midrule
1efs & heteroduplex & 818 & 2 & OL21/OL3, OL15/OL3, OL24/OL3, CHARMM36 & DNA:RNA hybrid, model 1 of PDB 1EFS; carries the term-by-term comparison \\
1BNA & DNA duplex & 758 & 2 & OL21, OL15, CHARMM36, OL24 & Drew--Dickerson dodecamer d(CGCGAATTCGCG), B-DNA duplex \\
d(ACGT) & terminal control & 126 & 1 & OL21, OL15, OL24 & DNA 4-mer, one of each base, free $5^\prime$-OH and $3^\prime$-OH \\
2KOC & RNA hairpin & 447 & 1 & OL3, CHARMM36 & UUCG tetraloop RNA hairpin, model 1 of PDB 2KOC \\
r(ACGU) & terminal control & 127 & 1 & OL3 & RNA 4-mer, one of each base, free $5^\prime$-OH and $3^\prime$-OH \\
1UBQ & protein, globular & 1,231 & 1 & ff14SB, ff19SB & Ubiquitin, $\beta$-grasp fold \\
5PTI & protein, globular & 892 & 1 & ff14SB, ff19SB & BPTI, three disulfide bonds \\
ACE-AGSV-NME & protein, peptide & 56 & 1 & ff14SB, ff19SB & Capped tetrapeptide, four residue kinds, built extended \\
Tax 11--19 & protein, peptide & 165 & 1 & ff14SB, ff19SB & HTLV-1 Tax nonamer LLFGYPVYV, capped, free in solution \\
1HHK & protein--peptide complex & 6,139 & 3 & ff14SB, ff19SB & HLA-A*02:01, $\beta_2$-microglobulin and the bound Tax nonamer \\
1F7Y & protein--nucleic complex & 3,321 & 2 & ff14SB/OL3, ff19SB/OL3 & Ribosomal protein S15 bound to a 16S rRNA fragment \\
\bottomrule
\end{tabular}

\end{table}

\begin{figure}[tbp]
\centering
\includegraphics[width=\linewidth]{systems-gallery.pdf}
\caption{The eleven systems, drawn from the coordinates the engine reads.
Nucleic acids and proteins are shown as cartoons, the peptides and the two
free-ended 4-mers as sticks with heteroatoms coloured by element; hydrogens are
present in every file and omitted from the drawing. Chain colour follows the
polymer, with the second protein chain of 1HHK ($\beta_2$-microglobulin) in a
lighter blue and its bound Tax nonamer as violet sticks.}
\label{fig:systems}
\end{figure}

\paragraph{1efs}
A heteroduplex of one DNA strand and one RNA strand, model~1 of Protein Data
Bank entry 1EFS\cite{pdb1efs,hantz2001hybrid}. Chain A is
\texttt{DG DA DG DA DG DG DA DA DG DA DG DA DA}, six DG at 33 atoms each and
seven DA at 32, giving 422 atoms; chain B is
\texttt{RU RU RC RU RC RU RU RC RC RU RC RU RC}, seven RU at 30 atoms and six
RC at 31, giving 396. Each residue carries a charge of $-1$, so the net charge
is $-26$ and the run uses \verb|\neutralize{off}|. It is the only structure
here that mixes DNA and RNA in one duplex, and that is why it carries the
term-by-term comparison: a parameter set correct for DNA and wrong for RNA is
invisible on a pure duplex of either, and this is the system on which F-03
became visible. The coordinates are taken as deposited apart from six built
atoms, a $5'$ phosphate on each strand, and four renaming or renumbering edits,
all listed in Supporting Information section~S7. The hydrogens are the
deposited NMR hydrogens; none was rebuilt, and nothing was minimised. All four
AMBER presets build the same padded system of 942 atoms in 30 residues, then
trim four terminal residues to the 818 atoms compared, so every residue that
enters the term-by-term comparison is an interior one. The CHARMM36 cell on
this system is a different molecule: the distributed XML has no patch that both
keeps a $5'$ phosphate and ends a chain (F-08), so the structure was rebuilt as
a $5'$-OH/$3'$-OH strand, 816 atoms, six phosphate atoms removed and four
hydroxyl hydrogens added, and its total is not comparable with the AMBER rows.

\paragraph{Capped tetrapeptide}
ACE-ALA-GLY-SER-VAL-NME, 56 atoms, with four residues of four different kinds
between the two caps. The four kinds are deliberate: each interior residue
produces one $(\phi,\psi)$ correction term, ff19SB carries a different map for
different residues, and a peptide with one interior residue would resolve its
single term through a fallback and return the same energy whether the map
selection were right or wrong. The coordinates come from \texttt{tleap} and are
not minimised: what \texttt{tleap} writes is an extended chain at
$\phi=\psi=180^{\circ}$ on every residue. Two consequences run through the
paper. The out-of-plane terms sit exactly at their minima, so both MOSAICS and
GROMACS report zero for them and the comparison on that term is vacuous
(F-11); and all four correction-map terms sit at $\phi=\psi=\pi$, the singular
point of the implemented force expression, which is why every energy comparison
passed while the analytic CMAP force was identically zero (F-21).

\paragraph{The free-ended 4-mers}
d(ACGT), 126 atoms, and r(ACGU), 127 atoms, one of each base with free $5'$-OH
and $3'$-OH ends. They isolate the terminal residues, which is where the residue
templates of two force fields differ, and they carry the force sweeps, the
CHARMM36 comparison and, through the terminal decks, AMBER comparisons against
sander and OpenMM. d(ACGT) carries a geometric defect at residue DG~3, finding
F-17, set out with its consequences in Section~\ref{sec:limits}.

\paragraph{The deposited structures}
1BNA, 1UBQ, 5PTI and 2KOC were fetched from the RCSB on 25 August 2026 and
prepared under one set of rules, recorded in
\texttt{evidence-v2/\_structures/BUILD-DECISIONS.md}. Only \texttt{ATOM}
records survive; every water, ion and ligand is dropped. Alternate location A
is kept. Every hydrogen and deuterium in the deposit is stripped and rebuilt by
the topology builder, so hydrogen placement is \texttt{tleap}'s for the AMBER
decks and the OpenMM template's for the CHARMM path, never the depositor's.
Protein termini are charged, NH$_3^+$ and COO$^-$, because a whole protein is
not a fragment; nucleic termini are $5'$-OH and $3'$-OH. 1BNA is the
Drew--Dickerson dodecamer d(CGCGAATTCGCG)\cite{pdb1bna}, both strands kept
because the duplex is self-complementary and one strand is not the dodecamer,
758 atoms, net charge $-22$; its terminal residues build as DC5 and DG3. 2KOC
is model~1 of the 20-model NMR ensemble of the 14-mer cUUCGg tetraloop
hairpin\cite{pdb2koc}, GGCACUUCGGUGCC, 447 atoms; its deposited $5'$-terminal
phosphate is removed so that the chain starts at a $5'$-OH, because a bare
$5'$ phosphate is the one terminus the CHARMM36 XML cannot end a chain on, and
keeping it would have bought a second system CHARMM36 cannot run. 1UBQ is
ubiquitin\cite{pdb1ubq}, chain A, 76 residues, 1,231 atoms, net charge 0; its
one histidine, HIS~68, is built as HIE, and the C-terminal tail deposited at
occupancy 0.25--0.45 is taken as deposited. 5PTI is bovine pancreatic trypsin
inhibitor from the joint X-ray and neutron refinement\cite{pdb5pti}, chain A,
58 residues, 892 atoms, net charge $+6$; its three disulfides, CYS~5--55,
14--38 and 30--51, are built as CYX and bonded in \texttt{tleap} exactly as the
deposited SSBOND records give them, and the 103 deuterium positions of the
neutron model are stripped with the hydrogens.

\paragraph{The peptide and the two complexes}
The Tax nonamer LLFGYPVYV, the HTLV-1 epitope presented by HLA-A*02:01, is
taken from the engine's own example of that peptide, capped with ACE and NME,
165 atoms, at the main-branch commit the run record names. 1HHK is the crystal
structure of HLA-A*02:01 with $\beta_2$-microglobulin and that
nonamer\cite{pdb1hhk}, three chains of 275, 100 and 9 residues, built with
\texttt{tleap} under ff14SB and ff19SB from the prepared coordinates, 6,139
atoms with 2,978 hydrogens added, net charge $-9$. 1F7Y is ribosomal protein
S15 bound to a fragment of 16S rRNA\cite{pdb1f7y}, chain A an 86-residue
protein with charged termini and chain B a 57-nucleotide RNA, 3,321 atoms, net
charge $-49$. It is the first system run against a deck that holds protein and
RNA atom types together, and those two decks live on an engine branch that has
not been merged, which Table~\ref{tab:matrix} records. Its reference
coordinates were written from the three-decimal PDB the engine reads rather
than from \texttt{tleap}'s restart file, because the restart differs from the
PDB by up to $5\times10^{-4}$~\AA{}, which is worth 0.66~\kcal{} of
Lennard-Jones energy on this system.

One rule governs every build. \emph{The structure MOSAICS reads is the
structure the reference topology is built from, never the reverse.} MOSAICS
reads the PDB and the reference engines read the topology, and the two formats
store positions to different numbers of decimals; the largest resulting
difference on the heteroduplex was 0.0005~\AA{}. On the peptide that moved the
short-range energy by 0.573~\kcal{}, because that term rises as the twelfth
power of distance, and it looked exactly like a broken force field (F-02). The
builders write the reference positions from the structure MOSAICS reads and
refuse to continue if the two disagree.

\subsection{Run conditions}
\label{sec:methods-conditions}

Every single-point run is in vacuum, at dielectric 1, with no cutoff, no
temperature coupling, no solvent and no move. No temperature enters any number
reported in Sections~\ref{sec:coverage} to~\ref{sec:forces}: nothing is
propagated, and the one proposal that is made reproduces the current structure,
so the energy difference it is judged on is zero. A comparison counts only if
the structure did not move, and the runs check that. MOSAICS prints each part
of a term next to that term's total, so the parts must add up. On 1efs with
OL21/OL3 they do so to within $5.5\times10^{-11}$~\kcal{}, the summation order
of a double-precision sum over 818 atoms; a larger discrepancy means the
structure moved.

\subsection{The electrostatic constant}
\label{sec:methods-coulomb}

The four programs do not share the constant that converts charge into energy,
and on 1efs the disagreement is worth more than everything else being measured.
It was measured, not looked up. One AMBER topology holding two point charges,
$+1\,e$ and $-1\,e$, exactly $1.0000000000$\,\AA{} apart, with Lennard-Jones
zeroed, in vacuum and with no cutoff, is handed to every engine; the
electrostatic energy an engine then prints \emph{is} its constant. The script
is \texttt{measure\_coulomb\_constant.py} in
\texttt{evidence/coulomb-constants/}, and the four values it returns are in
Table~\ref{tab:coulomb}. sander's is measured to four digits because it prints
four decimals; its exact value is $18.2223^2 = 332.052217$, because AMBER
stores topology charges pre-multiplied by 18.2223. MOSAICS's row is analytic
rather than measured: it is the product of two constants compiled into the
engine, $0.529177\times627.50921 = 332.06344122017$, and that unrounded product
is the number every correction is done with.

Before any comparison, the electrostatic term of the MOSAICS and sander columns
is scaled by the ratio of its own constant to a single target. That target is
not any engine's measured value but a number derived from the 2019 SI values of
$e$ and $N_{\mathrm{A}}$, the thermochemical calorie and the CODATA 2018
$\varepsilon_{0}$, given in Section~\ref{sec:constant}: $k = 332.0637133$
kcal\,\AA\,mol$^{-1}$\,$e^{-2}$. The OpenMM and GROMACS columns are left as
printed: both print six decimals and to those six decimals both equal the
derived target. Every total in the 26-cell matrix and every electrostatic
residual in the term-by-term tables is on this constant.

\subsection{Force validation}
\label{sec:methods-forces}

The analytic force is compared with a central difference of the same energy,
atom by atom and component by component, for each term separately and for the
total. The step is swept over ten values from $10^{-2}$ to
$3\times10^{-7}$\,Bohr, and what is recorded at each step is the largest
$|$numerical $-$ analytic$|$ over every atom and component, in Hartree per
Bohr. Sixty-one such sweeps are committed in \texttt{evidence/forces/}:
thirty-seven against Cartesian coordinates on 1efs, the two 4-mers and the
peptide, and twenty-four against a staging coordinate, the fourth staging
system being a capped peptide fixture whose NME cap is not a physical amide and
whose energy is therefore recorded, not validated (F-18).

The measurement is the slope, not the residual. A central difference is wrong
at order step$^2$, so a force that is the derivative of the energy makes the
residual fall as the square of the step until round-off in the energy takes
over, while a force wrong by a fixed amount stops falling at that amount. The
slope is fitted by least squares on $\log_{10}$(residual) against
$\log_{10}$(step) over a window declared in advance, $10^{-2}$ to
$3\times10^{-4}$\,Bohr, four of the ten steps and spanned by every sweep.
Fixing the window matters, because a fit over rows selected for already
falling as the square would test nothing.

\subsection{The cost of one Monte Carlo step}
\label{sec:methods-cost}

A single run is setup, parse, one initial energy, $N$ steps and output. At
$N=1$ the constant part is 99.8\,\% of the wall clock, so timing one run and
calling it the cost of one evaluation would overstate it by more than two
orders of magnitude. The constant was measured instead of assumed. The
committed 1efs OL21/OL3 input was run unmodified except for
\verb|\total_step_mc| at $N = 0, 1, 5, 10, 25, 50$ and $100$, five repeats
each, and $t(N) = a + bN$ was fitted by least squares on the medians; $b$ is
the marginal cost of a step and $a$ is everything that happens once. Medians
rather than means because this is a laptop: the $N=0$ point has one 866~ms
outlier against a 442~ms floor. The build was the serial target at commit
\texttt{eccdf74}, Apple clang 21.0.0 at \texttt{-O3}, one process on one core
of an Apple M5, and the timed run is the committed run: its initial total
energy, 2526.184857387672~\kcal{}, is byte-identical to the committed
single-point output. Acceptance over the timed steps was 59 of 100. A rejected
step still costs a full energy evaluation, since the evaluation is what decides
it, so the slope is an average over a stable accept/reject mix and not a cost
conditional on acceptance.

\subsection{\studio{}}
\label{sec:methods-studio}

\studio{} is a web page. Behind it is a Python package (Python 3.10 or later,
MIT licence) whose pipeline is headless: four functions carry every operation,
\path{sequence_to_structure}, \path{structure_to_minimised},
\path{build_variant} and \path{compare_structures}, and the page calls those
functions through Pyodide. The engine is the MOSAICS executable compiled
to WebAssembly and driven through a small JavaScript bridge: the pipeline
writes the input file, the residue topology, the deck and the structure into a
run folder on the in-tab file system, runs the engine on it, follows its
progress from the lines it prints and reads back every file it writes.

\paragraph{The guided workflow}
The page presents four steps over that pipeline. \emph{Start} takes a sequence (DNA, RNA or a peptide, a pasted or
dropped FASTA record) or a coordinate file, and builds nucleic acids on ideal
A- or B-form geometry as single strands, duplexes or DNA$\cdot$RNA hybrids;
peptides are built on an ideal backbone with side chains and pH-7 protonation
from PDB2PQR's residue definitions\cite{dolinsky2004pdb2pqr,jurrus2018apbs}
and optional ACE/NME caps; proteins are prepared from coordinates, not built.
\emph{Prepare} chooses the chains to keep, leaves ligands and waters out, adds
hydrogens to a deposited protein, and chooses the force field from every
compatible profile with the reason the default was chosen; thirteen profiles
are declared, twelve over the AMBER and CHARMM36 decks of
Section~\ref{sec:methods-parameters} and one knowledge-based three-point
protein model, and each names its six files. One validation gate
sits on this step: whether the ordinary bonded and non-bonded energies for
every residue type present reproduce a committed reference. \emph{Run} offers
fifteen protocols, each with the deck it comes from, the work it is cited to
and a link to the engine's manual, and exposes every parameter the engine reads
for that kind of calculation, grouped as the manual groups them; the deck the
engine will read is one click away. \emph{Understand} shows the energy along
the run, the RMSD to the starting structure along the trajectory, energy
against RMSD with one point per frame, acceptance tallies, the energy table, a
residue-by-residue comparison of before and after, the molecular graph export
(NumPy, JSON, PyTorch Geometric and DGL formats
\cite{fey2019pyg,wang2019dgl}, with no learned embeddings), and the record.
The three-dimensional view is 3Dmol.js\cite{rego2015threedmol}, and a viewer
plays the trajectory frame by frame.

\paragraph{The record}
Every run writes a record naming the Studio version, the engine, the SHA-256 of
the engine executable, the SHA-256 of each force-field file that was staged,
the preset and every parameter it resolved, the inputs and the outputs by
digest, and the start and end times. The record and every file it names
can be downloaded from the page, and the package's verify routine recomputes
every digest and reports whether the recorded files are present and unchanged.
The engine itself already writes every parameter it resolved to
\texttt{sim\_param.out}; what the record adds is the binding of those
parameters to the engine build, the deck files and the inputs in one document
that can be re-checked.

\paragraph{The engine in the browser}
The engine was compiled to WebAssembly with Emscripten
6.0.9\cite{zakai2011emscripten} from the serial target of the engine's main
branch at commit \texttt{73f598a} (7 September 2026). The 151 objects compile
with no source changes; the link needs a 32~MB stack, because the input parser
overruns Emscripten's 64~KB default at once, and one guard on the single
\texttt{popen("date")} call, made in the scratch copy used for the build and
not in the engine repository. The result, \texttt{mosaics.wasm}, is 713,184
bytes. The Studio's own pipeline runs beside it under Pyodide
0.27.7\cite{pyodide-software} in a Web Worker; one Python shim registers the
WebAssembly build as the engine and replaces the one function that would have
launched a native executable, and everything else is the package's own code. The
distribution is about 8.5~MB of static files (the page, the engine, the Python
archive, one archive per force-field profile fetched only when that profile is
chosen, the manuals and a version manifest); a first visit downloads about
25~MB including Pyodide and NumPy from a content delivery network, cached
afterwards. Every URL the page fetches carries the build identifier, so a
redeploy is never served from a browser's cache. Nothing is uploaded and no
server computes: any static host serves it, and the copy reported here is on a
personal web page at the University of Oxford. The engine is single-threaded
in the browser, a run cannot be interrupted from outside, and everything lives
in the tab's memory until it is downloaded.

\paragraph{Tests}
The package carries 84 test files with 2,198 test functions (counted on 8
September 2026), run on every push on three operating systems across Python
3.10 to 3.13, and the browser build is checked end to end by a headless test
that drives the whole pipeline under Node with the same Python archive and the
same bridge the page uses, and fails unless the documented case study
reproduces (Section~\ref{sec:studio-results}).
